\documentclass {article}
\usepackage{graphicx}
\usepackage{amsmath}
\usepackage{amssymb}
\renewcommand\vec{\mathbf}
\let\OldS\nabla
\renewcommand{\nabla}{\boldsymbol{\OldS}}

\let\OldHat\hat
\renewcommand{\hat}[1]{\OldHat{\mathbf{#1}}}


\begin{document}
\thispagestyle{empty}% empty header/footer

\begin{center}\strut
	\bfseries\Huge
	What if  Mass Eats the Vacuum? 

--------------------------------


	Mass as a Sink of Zero-Point Radiation: A Classical Alternative to General Relativity
\end{center}


\centerline{David Plotz}


\begin{center}\strut
May 5th, 2025
\end{center}
\vfill

	
\newpage
\begin{abstract}
	This paper introduces a classical field theory of gravity built on modified Maxwell equations and a central physical hypothesis: that mass absorbs zero-point radiation, inducing a net inward momentum flow. Drawing from stochastic electrodynamics (SED), we postulate a Lorentz-invariant radiation field present throughout space, and hypothesize that mass acts as a sink for this field up to a frequency cutoff \( \omega_{\text{max}} = m / \hbar \).  We construct a consistent model of a massive particle, demonstrate energy and momentum conservation, and derive an emergent Dirac equation via a spinor transformation. Classical tests of general relativity—Mercury's precession, gravitational lensing, and gravitational redshift—are reproduced to leading order. This approach provides a novel, Lorentz-invariant, classical alternative to quantum gravity and general relativity, and lays a foundation for further exploration.
\end{abstract}
\newpage

\section{Introduction}

Physics today stands at an impasse. The Standard Model and general relativity have provided extraordinary y insight into the workings of the universe, but their domains remain disjoint: quantum mechanics governs the small, general relativity the large. Despite nearly a century of effort, these two frameworks remain fundamentally incompatible.

Efforts to unify them—most notably string theory—have produced elegant mathematics but no testable predictions. The field of fundamental physics, as argued by Lee Smolin and others, has stagnated, dominated by a narrow set of ideas and institutional inertia. We are still without a coherent, testable theory of gravity that naturally connects to quantum phenomena.

This paper proposes a radically different approach. Rather than quantizing gravity or geometrizing quantum mechanics, we hypothesize a classical field-theoretic mechanism by which gravitational behavior emerges from radiation momentum transfer. The core idea is that mass acts as a \emph{sink} for zero-point radiation, inducing a directional flow of momentum in the surrounding field. 

The theory is formulated in terms of modified Maxwell-like equations, inspired by Heaviside’s gravitoelectromagnetism and enriched by the stochastic background field ideas of Boyer and the stochastic electrodynamics (SED) community. Throughout the paper, we maintain a classical framework: no quantization, no operator formalism, no probabilistic interpretation — yet key features of quantum and gravitational behavior emerge as natural consequences of the dynamics.

We offer an outline of the remaining paper: 

2. A Physical Picture: Radiation Sink Intuition

3. Maxwell-Style Equations for Gravity

4. Stochastic Electrodynamics

5. Coupling Mass to the Field: Dual Cutoffs

6. Field Energy and Momentum Balance: A Mechanism for Stability

7. Classical Tests of Gravity

8. Duality and the Dirac Equation

9. Discussion and Outlook

\section{A Physical Picture: Radiation Sink Intuition}

Before diving into equations, we begin with an image. Picture a warm iron ball, fresh from a furnace, glowing red as it radiates heat. Assume it is electromagnetically neutral and a perfect conductor.

Now place this ball in the vacuum of deep space. As it cools, it emits infrared photons, gradually losing energy to its surroundings. The glow fades, and the ball approaches absolute zero.

But even at zero Kelvin, the universe is not truly empty. Both classical and quantum considerations point to a persistent electromagnetic background: the zero-point field (ZPF). A fuller treatment of this will be offered in the section Stochastic Electrodynamics. For now, imagine a faint hiss of zero-point radiation filling even the coldest vacuum. The iron ball, though no longer glowing, continues to absorb from this ambient field.

Our proposal is this: even at absolute zero, mass interacts with the ZPF. The iron ball draws in energy—not passively, but by creating an inward flow of momentum from the field. This absorption forms a radial momentum flux toward the mass. The cold ball doesn't just sit inert—it sinks radiation from the ZPF.


This is the metaphorical heart of the theory. Just as electric charges source field lines, and magnetic poles source curls, we propose that mass acts as a sink in the momentum field of zero-point radiation. The rest of the paper builds from this intuition.

\subsection{Main Hypothesis and 2 assumptions}


This image above anchors the theory: 


if mass always interacts with the zero-point field, then this interaction must have physical consequences. We propose it manifests as an inward-pointing momentum density:

\[
\boxed{
	\vec{p}(r) = -\hat{r} \frac{G m \hbar}{r^2}
}
\]


This is not yet derived—it is proposed. But it prompts key questions: Can such a flow obey energy and momentum conservation? Can it stabilize mass against collapse? Can it account for gravitational phenomena?

Here, we also simply state the hypothesis and show its consequences, we hope in future drafts of this paper or in subsequent work that we will instead make this hypothesis simply a result of the 2 assumptions we will make in the next two sections and alluded to earlier..

In the following section, we begin quantifying the limits: How small can mass become before gravitational collapse takes over? Up to what frequency does it couple to the ZPF?


\section{Maxwell-Like Equations for Gravity}

Our gravitational field theory begins with a structural analogy: just as electromagnetism is governed by Maxwell's equations, we propose that gravity obeys a similar set of vector field equations.  These equations follow the form of Maxwell's equations in electromagnetism, but with mass and mass current acting as sources rather than charge and electric current. In natural units (\( c = 1 \)):

\[
\nabla \times \vec{B} - \partial_t \vec{E} = -4 \pi G \vec{J}, \quad 
\nabla \times \vec{E} + \partial_t \vec{B} = 0
\]
\[
\nabla \cdot \vec{E} = -4 \pi G \rho, \quad 
\nabla \cdot \vec{B} = 0
\]


These equations describe a vector field sourced by mass density \( \rho \) and mass current \( \vec{J} \), with \( G \) playing the role of the coupling constant. Crucially, this is a **classical** theory: there is no operator formalism or quantum interpretation. Fields are continuous, real, and obey Lorentz invariance.

\subsection{Irrotational and Solenoidal Decomposition}

Any vector field can be split into irrotational (curl-free) and solenoidal (divergence-free) components:

\[
\vec{V} = \vec{V}_I + \vec{V}_S, \quad 
\nabla \cdot \vec{V}_S = 0, \quad 
\nabla \times \vec{V}_I = 0
\]

Applying this to our gravitational fields yields two decoupled subsystems:

\vspace{0.5em}
\noindent\textbf{Static Fields (Irrotational):}
\[
\nabla \cdot \vec{E}_I = -4 \pi G \rho, \quad 
\partial_t \vec{E}_I = 4 \pi G \vec{J}_I, \quad 
\nabla \cdot \vec{J}_I + \partial_t \rho = 0
\]

\noindent\textbf{Radiation Fields (Solenoidal):}
\[
\nabla \times \vec{B}_S - \partial_t \vec{E}_S = 4 \pi G \vec{J}_S, \quad 
\nabla \times \vec{E}_S + \partial_t \vec{B}_S = 0
\]

This decomposition is not just convenient—it is physically meaningful. This decomposition is more than formal. It physically separates the **static gravitational attraction**—which reduces to Newton’s law in the appropriate limit—from the **propagating radiation** field, which transfers energy and momentum.

Because these sectors are decoupled (in the rest frame of the mass), we can analyze them independently. This gives us flexibility: we can examine the radiation structure in full detail without disturbing the familiar gravitational attraction.

\subsection{Recovery of Newtonian Gravity}

In the case of a static point mass \( m \) located at the origin, we set:

\[
\rho(\vec{x}) = m \delta^3(\vec{x}), \quad \vec{J} = 0
\]

From the irrotational equation \( \nabla \cdot \vec{E}_I = -4 \pi G \rho \), the solution is the familiar inverse-square law:

\[
\vec{E}_I(\vec{x}) = - \frac{G m}{r^2} \hat{r}
\]

This recovers Newtonian gravity exactly in the weak-field limit. No spacetime curvature is required—only field lines and momentum flow.

\subsection{Why a Magnetic-Like Field?}

Some may be surprised by the presence of \( \vec{B} \), a magnetic-like field. But this emerges naturally: just as a moving charge generates a magnetic field, a moving mass generates a “gravitomagnetic” field. This is well known from general relativity in the form of frame-dragging and gravitomagnetism, and our equations make it explicit. Static objects generate only \( \vec{E}_I \); moving objects induce \( \vec{B}_S \) and radiate.

\subsection{Radiation from the Solenoidal Field}

Outside the mass, the solenoidal field equations in vacuum are:

\[
\nabla \times \vec{E}_S + \partial_t \vec{B}_S = 0, \quad 
\nabla \times \vec{B}_S - \partial_t \vec{E}_S = 0
\]

Assuming time dependence \( e^{-i \omega t} \), the general solution is expanded in vector spherical harmonics:

\[
\vec{B}_S = \sum_{l,m} \left[
a_E(l,m) f_l(kr) \vec{X}_{lm}
- \frac{i}{k} a_M(l,m) \nabla \times \left( g_l(kr) \vec{X}_{lm} \right)
\right]
\]

\[
\vec{E}_S = \sum_{l,m} \left[
\frac{i}{k} a_E(l,m) \nabla \times \left( f_l(kr) \vec{X}_{lm} \right)
+ a_M(l,m) g_l(kr) \vec{X}_{lm}
\right]
\]

Here, \( f_l \) and \( g_l \) are spherical Hankel functions. Since we model mass as a **sink** for radiation, we include only incoming waves.

The energy flux into the particle is computed via the Poynting vector:

\[
P = \int \nabla \cdot \left( \frac{1}{8\pi} \vec{E} \times \vec{B}^* \right) \, d^3x 
= R^2 \int \hat{r} \cdot \vec{S} \, d\Omega \Big|_{r=R}
\]

Using mode orthogonality and assuming \( |a_E|^2 = |a_M|^2 \), the total absorbed power becomes:

\[
P = \sum_{l,m} \frac{|a_E|^2}{4 \pi k^2}
\]

The associated energy and momentum densities in the radiation zone also follow cleanly. The total field momentum is directed radially inward:

\[
\vec{P}_{\text{field}} = -\hat{r} \sum_{l,m} \int_V \frac{|a_E|^2 |\vec{X}_{lm}|^2}{4\pi k^2 r^2} \, d^3x
\]



The total field momentum:

\[
\vec{P}_{\text{field}} = \frac{1}{8\pi} \int_V \vec{E} \times \vec{B}^* \, d^3x
= -\hat{r} \sum_{l,m} \int_V \frac{|a_E|^2 |\vec{X}_{lm}|^2}{4\pi k^2 r^2} \, d^3x
\]

The specific coefficients \( |a_E|^2 \) will be fixed by boundary conditions in the section on collapse and stability. But, since we've already speeled out the hypothesis then we know we will eventually get to
:

\[
\boxed{
	\vec{p}(r) = -\hat{r} \frac{G m \hbar}{r^2}
}
\]

This is not a free assumption. In later sections, we will show how this momentum field arises from boundary conditions at the mass’s surface, and how the coefficient (including \( \hbar \)) follows from coupling to the zero-point background up to a frequency cutoff \( \omega_{\text{max}} = m / \hbar \).


\subsection{What Comes Next}

With the Maxwell-style field equations and mode structure in hand, we are now ready to couple this framework to the zero-point radiation background. In doing so, we move from a kinematic field theory to a dynamical one—where mass absorbs, redirects, and stabilizes through its interaction with the ambient field.


\section{Stochastic Electrodynamics (SED)}

To complete our framework, we now introduce the physical source of the radiation field: the zero-point background. Rather than assuming empty space, we posit—following the ideas of stochastic electrodynamics (SED)—that the vacuum is filled with a classical, Lorentz-invariant radiation field even at absolute zero. This field supplies the fluctuations from which mass sinks momentum.

\subsection{The Zero-Point Field}

In SED, the vacuum is not silent. It carries a random, omnipresent electromagnetic field with spectral energy density:

\[
\rho(\omega) = \frac{1}{2} \hbar \omega
\]

This spectrum is Lorentz-invariant and corresponds to the ground state energy per mode, in analogy to the quantum harmonic oscillator. But unlike quantum theory, which treats this as an abstract baseline, SED models it as a real classical field.

The field is stochastic and Gaussian, defined by:

\[
\langle W \rangle = 0, \qquad 
\langle W_i(\vec{k}) W_j^*(\vec{k}') \rangle = \delta_{ij} \delta^3(\vec{k} - \vec{k}')
\]

The electric and magnetic fields are then:

\[
\vec{E}(\vec{x}, t) = \text{Re} \sum_{\lambda=1}^2 \int d^3k \, 
\hat{\epsilon}(\vec{k}, \lambda) 
\sqrt{\frac{\hbar \omega}{2\pi^2}} 
W_\lambda(\vec{k}) e^{i(\omega t - \vec{k} \cdot \vec{x})}
\]

\[
\vec{B}(\vec{x}, t) = \text{Re} \sum_{\lambda=1}^2 \int d^3k \, 
\frac{\vec{k} \times \hat{\epsilon}(\vec{k}, \lambda)}{k}
\sqrt{\frac{\hbar \omega}{2\pi^2}} 
W_\lambda(\vec{k}) e^{i(\omega t - \vec{k} \cdot \vec{x})}
\]

This background field is homogeneous, isotropic, and defined over all frequencies up to some high-frequency cutoff. We assume that each mode contributes an energy \( \tfrac{1}{2} \hbar \omega \), and that the full ensemble of modes is statistically stationary and invariant under Lorentz transformations.

\subsection{Why This Matters}

This field is not added as an afterthought—it is the heartbeat of the theory. Without it, mass would generate fields but receive nothing in return. With it, a dynamic equilibrium becomes possible. Mass interacts with the zero-point field, selectively absorbing momentum from it. This interaction, we propose, is the **source of gravity**.

In SED, zero-point radiation explains how classical systems can remain stable, even while radiating energy. It has been shown to reproduce:

- The Planck blackbody spectrum (Boyer, 1969)
- Harmonic oscillator ground-state energy \( \tfrac{1}{2} \hbar \omega \)
- Casimir and van der Waals forces
- Some features of the hydrogen atom ground state
- The Unruh effect

Each of these suggests that the vacuum is not passive. It is a sea of fluctuating modes that constantly jostle matter—even at \( T = 0 \).


\subsection*{Transition}

Now that the preliminaries are in place—namely, Maxwell’s equations (as adapted by Heaviside) and Stochastic Electrodynamics (as introduced by Boyer)—we can turn to the mathematics that is novel in this theory.

These two ingredients are not original to this work, but together they set the stage. What follows is where the new structure begins to emerge.

\section{Coupling Mass to the Field: Dual Cutoffs}

A central feature of this theory is that mass interacts with the zero-point field only over a limited domain, defined by two natural cutoffs: one in space, and one in frequency. These cutoffs are not imposed arbitrarily, but emerge from physical consistency. The spatial cutoff prevents gravitational collapse, while the frequency cutoff ensures that absorption from the zero-point field remains finite.

\subsection*{Interaction Hypothesis}

We propose that a mass \( m \) couples to the zero-point field up to a maximum frequency 
\[
\boxed{\omega_{\text{max}} = \frac{m}{\hbar}}
\]
and within a radial domain extending at least to 
\[
\boxed{r \gtrsim Gm}
\]

This assumption is not yet derived, but it will prove physically consistent. In later sections, we show how this interaction yields a radial inward momentum density:
\[
\boxed{\vec{p}(r) = -\hat{r} \frac{G m \hbar}{r^2}}
\]
and how the constant in this expression emerges from the field equations and boundary conditions.

\subsubsection*{Lower Bound: Minimum Radius via Gravitational Self-Energy}

From classical considerations, any spherically symmetric mass distribution has a gravitational self-energy:
\[
U = -\frac{1}{2} \int \frac{G m \rho(\vec{x})}{r} \, d^3x
\]

For a spherical shell of radius \( R \), this yields:
\[
R = \frac{G m}{2}
\]

Other geometries differ only by numerical prefactors. In general, we require:
\[
R \gtrsim G m
\]

We interpret this as a \textbf{minimum radius} for any stable mass. Unlike the Schwarzschild radius (which assumes collapse), this lower bound emerges as a stability criterion under constant zero-point radiation pressure.

\subsubsection*{Upper Bound: Cutoff Frequency via Energy Matching}

Conversely, mass cannot absorb arbitrarily high-frequency radiation. A natural upper limit arises from energy conservation:
\[
U = m = \hbar \omega_{\text{max}} \quad \Rightarrow \quad \omega_{\text{max}} = \frac{m}{\hbar}
\]

This sets the upper bound for momentum transfer and defines the domain of coupling between mass and field. While this frequency is extraordinarily high for macroscopic objects, it plays a crucial theoretical role in regularizing energy and momentum integrals.

\vspace{0.5em}
Together, these dual cutoffs—one in radius, one in frequency—define a resolution scale for gravity. They constrain how mass interacts with the field and suggest a breakdown of classical behavior beyond these limits, potentially pointing toward quantum or nonlinear extensions of the theory.



\section{Field Energy and Momentum Balance: A Mechanism for Stability}

\subsection{Stability from Field Energy}

We now evaluate whether the inward radiation flow, hypothesized in Section 2, can provide a stabilizing mechanism for a massive particle. Specifically, we ask: can the absorbed energy from the zero-point field counteract the gravitational self-energy that would otherwise lead to collapse?

The total field energy is given by:

\[
U = \frac{1}{16\pi} \int \left( |\vec{E}|^2 + |\vec{B}|^2 \right) \, d^3x
\]

From earlier results, the energy density near a massive particle absorbing zero-point radiation falls off as \(1/r^2\). Including the cutoff frequency \(\omega_{\text{max}} = m/\hbar\), we have:

\[
u(r) = \int_0^{m/\hbar} \frac{\alpha}{r^2} \, d\omega = \frac{\alpha m}{\hbar r^2}
\]

Integrating over all space, assuming spherical symmetry:

\[
U_{\text{rad}} = \int_0^\infty u(r) \cdot 4\pi r^2 \, dr = \int_0^\infty \frac{4\pi \alpha m}{\hbar} \, dr
\]

This diverges unless we impose a lower bound, namely the minimum radius of the particle \(R = Gm\). Truncating the integral from \(R\) to \(\infty\), we obtain:

\[
U_{\text{rad}} = \frac{4\pi \alpha m}{\hbar} \int_R^\infty dr = \infty
\]

To regularize this, we instead take the volume integral up to some finite domain or assume that most of the energy is absorbed near \(R\). The key is not the total energy, but the **balance** between the field energy and the self-gravitational energy of the particle:

\[
U_{\text{grav}} = -\frac{G m^2}{R}
\]

\[
U_{\text{rad}} = \frac{4\pi \alpha m R}{\hbar}
\]

Balancing these energies (assuming \(U_{\text{total}} = 0\)) gives:

\[
-\frac{G m^2}{R} + \frac{4\pi \alpha m R}{\hbar} = 0 \quad \Rightarrow \quad \alpha = \frac{G m \hbar}{4 \pi R^2}
\]

Substituting back, we find the energy density is:

\[
u(r) = \frac{G m \hbar}{4\pi R^2 r^2}
\]

This shows that the particle's own radiation environment acts to resist collapse, provided the momentum flow structure is accepted. The directionality of \(\vec{p} \propto -\hat{r}\) also implies an inward pressure that balances gravitational pull.

\subsection{Momentum Conservation and “Ohm’s Law”}

We now turn to momentum conservation. The total momentum exchange in the system should obey:

\[
\int_V \left( \rho \vec{E}_I + \vec{J}_S \times \vec{B}_S \right) \, d^3x = 0
\]

This relation balances the impulse from the irrotational field (associated with gravitational mass) against the magnetic impulse from the solenoidal radiation field.

Assuming the solenoidal current obeys a linear response:

\[
\vec{J}_S = \sigma \vec{E}_S
\]

and harmonic time dependence \(e^{-i \omega t}\), the field momentum becomes:

\[
-\hat{r} \int \left( m \delta(\vec{x}) \cdot \frac{G m}{r^2} - \sigma \cdot \frac{G m \hbar}{r^2} \right) \, d^3x = 0
\]

Solving gives:

\[
\sigma = \frac{m}{\hbar} \delta(\vec{x})
\]

This expression defines a delta-function conductivity concentrated at the origin — an effective Ohm's law in which the mass acts as a point-like, perfect absorber for incoming radiation. Including conventional normalization gives:

\[
\sigma = \frac{m}{4\pi \hbar} \delta(\vec{x})
\]

This result is not imposed, but emerges from requiring consistency with the momentum flux model.

\subsection{Summary}

Together, the field energy and momentum conservation arguments demonstrate that the hypothesized momentum inflow can act as a self-consistent stabilizing structure. A massive particle that sinks zero-point radiation does not collapse — not due to internal structure, but due to interaction with the vacuum field itself.

	
\section{Classical Tests of Gravity}

This section demonstrates that the present theory reproduces three key experimental predictions of general relativity: the precession of Mercury's perihelion, the deflection of light near a massive object, and gravitational redshift. Each is derived from the momentum-transfer hypothesis introduced earlier.

\subsection{Precession of the Perihelion of Mercury}

In electrodynamics, a static Coulomb field transforms under Lorentz boosts into a configuration with both electric and magnetic components. Specifically, the fields of a charge moving at velocity \( \vec{v} \) are:

\[
\vec{E}' = \gamma \vec{E} - \frac{\gamma^2}{\gamma + 1} \vec{v} (\vec{v} \cdot \vec{E}), \quad
\vec{B}' = - \gamma \vec{v} \times \vec{E}
\]

where \( \gamma = 1 / \sqrt{1 - v^2} \).

By analogy, we assume the same Lorentz transformation laws apply to gravity. We define the gravitational field:

\[
\vec{E}_g = \vec{F} / m = - \frac{G M}{r^3} \vec{r}
\]

Then, under a Lorentz boost:

\[
\vec{E}_g' = \gamma \vec{E}_g - \frac{\gamma^2}{\gamma + 1} \vec{v} (\vec{v} \cdot \vec{E}_g), \quad
\vec{B}_g' = -\gamma \vec{v} \times \vec{E}_g
\]

In the low-velocity limit \( v \ll 1 \), we find:

\[
\vec{B}_g = - \vec{v} \times \vec{r} \cdot \frac{GM}{r^3}
\]

This introduces a "gravitomagnetic" interaction analogous to magnetism in electrodynamics.

We now consider Kepler’s problem. The Newtonian effective potential is:

\[
U_{\text{Newt}} = \frac{1}{2} m \left( \frac{dr}{dt} \right)^2 - \frac{GMm}{r} + \frac{l^2}{2 m r^2}
\]

where \( l = m |\vec{r} \times \vec{v}| \) is the orbital angular momentum.

A moving charge induces a magnetic moment \( \vec{m} = \frac{1}{2} e (\vec{r} \times \vec{v}) \). Replacing \( e \rightarrow m \), we define a gravitational magnetic moment:

\[
\vec{m}_g = \frac{1}{2} m (\vec{r} \times \vec{v})
\]

The potential energy of this moment in a magnetic field is:

\[
U = -\vec{m}_g \cdot \vec{B}_g = -\frac{GM m}{2} |\vec{r} \times \vec{v}|^2 / r^3
\]

This adds a correction to the potential energy. Including the reciprocal effect (from the motion of the central mass), we double this term and find the corrected potential:

\[
U_{\text{Einstein}} = \frac{1}{2} m \left( \frac{dr}{dt} \right)^2 - \frac{GMm}{r}
+ \frac{l^2}{2 m r^2} - \frac{GM l^2}{m r^3}
\]

This is the same correction obtained in general relativity to first order. As Carroll notes, general relativity produces this result exactly; in our theory, it emerges as a first-order approximation. Higher-order terms may offer a route to testing deviations.

\subsection{Deflection of Light}

We now consider the bending of light near a massive object. Our central hypothesis is that mass creates a momentum-transfer field:

\[
\vec{P}(r) = - \hat{r} \frac{G m \hbar}{r^2}
\]

A photon traveling through this field receives a cumulative momentum shift:

\[
\Delta \vec{p} = \int_{\text{path}} \vec{P}(r) \, dl
\]

Assuming a straight-line path passing at minimum distance \( b \), we integrate:

\[
\Delta \vec{p} = \int_{-\infty}^{\infty} - \frac{G m \hbar}{r^3} \vec{r} \, w \, dz
\]

The parallel momentum shifts cancel symmetrically. The transverse component yields:

\[
\Delta p_y = \hbar w \cdot \frac{2 G m}{b}
\]

The angle of deflection is then:

\[
\theta \approx \frac{\Delta p_y}{p_z} = \frac{2 G m}{b} \quad \Rightarrow \quad \alpha = \frac{4 G m}{b}
\]

This matches the deflection angle predicted by general relativity.

\subsection{Gravitational Redshift}

Finally, we consider the redshift of light climbing out of a gravitational well. As before, a photon gains or loses momentum depending on its path through the momentum field.

Using:

\[
\Delta p = \int_{\infty}^{R} \vec{P} \cdot \hat{r} \cdot w \, dr = \int_{\infty}^{R} - \frac{G m \hbar}{r^2} w \, dr = \hbar w \cdot \frac{G m}{R}
\]

The resulting frequency shift is:

\[
\hbar w' = \hbar w + \Delta p = \hbar w \left( 1 + \frac{G m}{R} \right)
\]

which, to first order, matches the prediction of general relativity for gravitational redshift.

		
\section{Duality and the Dirac Equation}

\subsection{Duality Transformations}

As discussed in Jackson (Section 6.11), the Maxwell equations can be extended to include both electric and magnetic charges:

\[
\nabla \times \vec{B} - \partial_t \vec{E} = \vec{J}_e, \quad
\nabla \times \vec{E} + \partial_t \vec{B} = -\vec{J}_m
\]

\[
\nabla \cdot \vec{E} = \rho_e, \quad
\nabla \cdot \vec{B} = \rho_m
\]

The magnetic densities are assumed to obey the same continuity equation as their electric counterparts.

Now consider the following duality transformation:

\[
\vec{E} = \vec{E}' \cos \xi + \vec{B}' \sin \xi, \quad
\vec{B} = \vec{B}' \cos \xi - \vec{E}' \sin \xi
\]

\[
\rho_e = \rho_e' \cos \xi + \rho_m' \sin \xi, \quad
\rho_m = \rho_m' \cos \xi - \rho_e' \sin \xi
\]

\[
\vec{J}_e = \vec{J}_e' \cos \xi + \vec{J}_m' \sin \xi, \quad
\vec{J}_m = \vec{J}_m' \cos \xi - \vec{J}_e' \sin \xi
\]

This transformation leaves the generalized Maxwell equations invariant. As Jackson notes, the distinction between electric and magnetic charge is, in this framework, a matter of convention. If all particles share the same magnetic-to-electric charge ratio, a duality rotation can eliminate magnetic charges entirely. This perspective opens the door for deeper reinterpretations of classical field theory.

\subsection{Spinor Representation of Field Operators}

We now examine a way to recast Maxwell’s equations using spinor-like objects.

Define the vector curl as a matrix operator:

\[
\nabla \times \vec{V} = g_i \partial_i \vec{V}
\]

with:

\[
g_1 = \begin{pmatrix}
0 & 0 & 0 \\
0 & 0 & -1 \\
0 & 1 & 0
\end{pmatrix}, \quad
g_2 = \begin{pmatrix}
0 & 0 & 1 \\
0 & 0 & 0 \\
-1 & 0 & 0
\end{pmatrix}, \quad
g_3 = \begin{pmatrix}
0 & -1 & 0 \\
1 & 0 & 0 \\
0 & 0 & 0
\end{pmatrix}
\]

Define the gamma matrices as (These matrices satisfy the algebra of infinitesimal rotations and generate the curl operator in 3D”):

\[
\gamma^0 = \begin{pmatrix}
-1 & 0 \\
0 & -1
\end{pmatrix}, \quad
\gamma^i = \begin{pmatrix}
0 & g_i \\
-g_i & 0
\end{pmatrix}
\]

These matrices are not the standard Dirac representation, but they serve to connect the structure of the field equations to the spinor formalism in a suggestive way. This could be further explored with tools from representation theory and Clifford algebras.

\subsection{Emergence of the Dirac Equation}

Recall from the section on momentum conservation that:

\[
\nabla \times \vec{B}_S - \partial_t \vec{E}_S = -\frac{m}{\hbar} \vec{E}_S \delta^3(\vec{x})
\]

Restricting attention to the surface of the particle (i.e., “on-shell”), we drop the delta function and write:

\[
\nabla \times \vec{B}_S - \partial_t \vec{E}_S = -\frac{m}{\hbar} \vec{E}_S
\]

Now, apply the duality transformation and write the solenoidal Maxwell equations in this form:

\[
-\partial_t \vec{E}_S + \nabla \times \vec{B}_S + \frac{m}{\hbar} \vec{E}_S = 0
\]

\[
-\partial_t \vec{B}_S - \nabla \times \vec{E}_S + \frac{m}{\hbar} \vec{B}_S = 0
\]

Using the spinor curl definition:

\[
-\partial_t \vec{E}_S + g_i \partial_i \vec{B}_S + \frac{m}{\hbar} \vec{E}_S = 0
\]

\[
-\partial_t \vec{B}_S - g_i \partial_i \vec{E}_S + \frac{m}{\hbar} \vec{B}_S = 0
\]

Now define a six-component field:

\[
\psi = \begin{pmatrix}
\vec{E}_S \\
\vec{B}_S
\end{pmatrix}
\]

We can combine the two equations above into a compact form:

\[
\left( \gamma^\mu \partial_\mu + \frac{m}{\hbar} \right) \psi = 0
\]

This has the exact form of the Dirac equation, as given in Sakurai’s Advanced Quantum Mechanics, Eq. 3.31.

This is a remarkable result: starting from classical field equations and a single added hypothesis, we arrive at a spinor field satisfying the Dirac equation—without invoking quantization or spin from the outset. It suggests that elements of quantum theory may emerge naturally from classical field dynamics under the right assumptions.



\section{Discussion and Outlook}

\subsection{A Field, Not a Force}

At first glance, the directional momentum flow described in this theory might resemble a gravitational force. But it is not. The core hypothesis postulates a momentum density field sourced by mass and arising from asymmetric absorption of zero-point radiation. The resulting momentum transfer acts on light but cannot be described by \( \vec{F} = m \vec{a} \), nor does it require acceleration. It is not a mechanical push—it is field interaction.

This distinction is critical. In this model, light does not bend due to spacetime curvature, but because its momentum evolves under the influence of a background field. Whether this counts as a “force” is secondary; mathematically, it is a field-driven momentum shift.

\subsection{Experimental Implications and Testability}

At low frequencies, this theory reproduces general relativity’s standard predictions to first order: perihelion precession, light deflection, and gravitational redshift. However, it introduces a natural high-frequency cutoff, \( \omega_{\text{max}} = m/\hbar \), that could lead to deviations at high energies.

This sets the stage for falsifiability. Could light near ultra-dense masses bend slightly differently than GR predicts at near-Planckian frequencies? Might gravitational lensing data or high-precision interferometry detect subtle anomalies? These are open but well-posed questions.

Importantly, the theory comes with a scale—where it may transition, fail, or reveal something deeper. That is a feature, not a bug.

\subsection{Scope and Limitations}

This is an explicitly classical theory. It makes no claims about entanglement, Bell violations, or nonlocality. It treats the zero-point field as real but not quantized. It does not attempt to reconstruct quantum field theory or gauge symmetry from scratch.

Instead, it isolates one question: can gravitational behavior be modeled as a classical momentum interaction with the vacuum field? Within that scope, the answer appears to be yes.

\subsection{Future Directions}

Several paths remain unexplored:

- **Nonlinear interactions**: How do multiple nearby masses interact? Do momentum fluxes interfere? Is the field superposable?

- **Thermodynamic interpretations**: Can this model be linked to entropic gravity or energy-entropy relations? The presence of a directional momentum field invites statistical reinterpretation.

- **Toward quantization**: Is there a consistent path to quantize this model? Can radiation absorption be discretized or coupled to standard QFT?

\subsection{Final Remarks}

This paper is a prototype. It closes no doors and makes no sweeping claims. Its contribution is a single, sharp hypothesis: that mass sinks radiation and generates a field of momentum flux. From this, classical gravity-like behavior—and elements of quantum structure—appear to follow.

This is offered not as a conclusion, but as a beginning. Collaborators welcome.




\end{document}